\documentclass{article}
\input{../../../../lib.sty}

\title{\texttt{fn make\_float\_to\_bigint\_threshold}}
\author{\MichaelShoemateAuthor}
\date{}

\begin{document}

\maketitle

\contrib
Proves soundness of the implementation of \rustdoc{measurements/noise_threshold/nature/fn}{make\_float\_to\_bigint\_threshold} in \asOfCommit{mod.rs}{f5bb719}.

\section{Hoare Triple}
\subsection*{Precondition}
\subsubsection*{Compiler-verified}

\begin{itemize}
    \item Generic \texttt{TK} implements trait \rustdoc{traits/trait}{Hashable}
    \item Generic \texttt{TV} implements trait \rustdoc{traits/trait}{Float}
    \item Const-generic \texttt{P} is of type \texttt{usize}
    \item Generic \texttt{QI} implements trait \rustdoc{traits/trait}{Number}
    \item Type \texttt{RBig} implements traits \texttt{TryFrom<TV>} and \texttt{TryFrom<QI>}.
        This is for fallible exact casting to rationals from floats in the function and input sensitivity in the privacy map.
    \item Type \texttt{i32} implements trait \rustdoc{traits/trait}{ExactIntCast}\texttt{<<T as FloatBits>::Bits>},
        This requirement means that the raw bits of \texttt{T} can be exactly cast to an \texttt{i32}.
\end{itemize}

\subsubsection*{Caller-verified}
None

\subsection*{Pseudocode}
\lstinputlisting[language=Python,firstline=2,escapechar=|]{./pseudocode/make_float_to_bigint_threshold.py}

\subsection*{Postcondition}
\begin{theorem}
    \validTransformation{(\texttt{input\_space, k, TK, TV, P, QI})}{\texttt{make\_float\_to\_bigint\_threshold}}
\end{theorem}

\begin{proof}
    In the definition of the function on line \ref{line:value-function},
    NaNs are excluded by the constructor, and clamping to finite bounds ensures that conversion to rationals succeeds.
    The precondition for \texttt{find\_nearest\_multiple\_of\_2k} is satisfied by line \ref{line:check-k},
    so \texttt{find\_nearest\_multiple\_of\_2k} is infallible.
    There are no other sources of error in the function,
    so the function cannot raise data-dependent errors.

    The function also always returns a hashmap with the same keys, and IBig values,
    meaning the output of the function is always a member of the output domain,
    as defined on line \ref{line:output-domain}.

    The stability argument breaks down into three parts:
    \begin{itemize}
        \item Clamping to finite bounds is 1-stable and conversion to rationals preserves the clamped values.
        \item Regard missing keys as having value zero. Rounding preserves zero and equal values remain equal,
        so the support of the difference cannot increase. Apply the postcondition of
        \rustdoc{measurements/noise/nature/float/utilities/fn}{get\_rounding\_distance}
        to the at most $\Delta_0$ differing coordinates for the $L_P$ bound,
        and to one coordinate for the $L_\infty$ bound.
        These allowances are added on lines \ref{line:lp-x-mul-2k} and \ref{line:li-x-mul-2k}, respectively.

        \item Rescaling the rounded values by $2^{-k}$ on line \ref{line:find-nearest} is $2^{-k}$-stable.
            This scaling is reflected on lines \ref{line:lp-x-mul-2k} and \ref{line:li-x-mul-2k},
            where the sensitivity is multiplied by $2^{-k}$.
    \end{itemize}

    The stability map checks that both $L_p$ and $L_\infty$ distances are nonnegative
    before applying the rounding allowance. Zero distances remain valid.

    For every pair of elements $x, x'$ in \texttt{input\_domain} and for every pair $(\din, \dout)$,
    where \din\ has the associated type for \texttt{input\_metric} and \dout\ has the associated type for \\ \texttt{output\_metric},
    if $x, x'$ are \din-close under \texttt{input\_metric}, $\texttt{stability\_map}(\din)$ does not raise an exception,
    and $\texttt{stability\_map}(\din) \leq \dout$,
    then $\function(x), \function(x')$ are $\dout$-close under \texttt{output\_metric}.
\end{proof}


\end{document}
